% Ad Atticum 14.13 (ad-atticum-14-13)
% Source: https://raw.githubusercontent.com/PerseusDL/canonical-latinLit/master/data/phi0474/phi057/phi0474.phi057.perseus-lat1.xml
% Fetched by scripts/fetch_latin.py

% Dateline (Perseus): Scr. in Cumano vi K. Mai. a. 710 (44).

\ciceroLetterOpener{CICERO ATTICO salutem}

\ciceroSection{1} septimo denique die litterae mihi redditae sunt quae erant a te xiii Kal. datae; quibus quaeris atque etiam me ipsum nescire arbitraris utrum magis tumulis prospectuque an ambulatione ἁλιτενεῖ delecter. est me hercule, ut dicis, utriusque loci tanta amoenitas ut dubitem utra anteponenda sit. a)ll' ou) daito\s e)phra/tou e)/rga me/mhlen, a)lla\ li/hn me/ga ph=ma, diotrefe/s, ei)soro/wntes dei/dimen: e)n doih=| de\ sawse/men h)\ a)pole/sqai.

\ciceroSection{2} quamvis enim tu magna et mihi iucunda scripseris de D. Bruti adventu ad suas legiones in quo spem maximam video, tamen si est bellum civile futurum (quod certe erit si Sextus in armis permanebit, quem permansurum esse certo scio), quid nobis faciendum sit ignoro. neque enim iam licebit quod Caesaris bello licuit neque huc neque illuc. quemcumque enim haec pars perditorum laetatum Caesaris morte putabit (laetitiam autem apertissime tulimus omnes), hunc in hostium numero habebit; quae res ad caedem maximam spectat. restat ut in castra Sexti aut, si forte, Bruti nos conferamus. res odiosa et aliena nostris aetatibus et incerto exitu belli , et nescio quo pacto tibi ego possum, mihi tu dicere, τέκνον ἐμόν, οὔ τοι δέδοται πολεμήια ἔργα, ἀλλὰ σύ γ' ἱμερόεντα μετέρχεο ἔργα λόγοιο

\ciceroSection{3} sed haec fors viderit, ea quae talibus in rebus plus quam ratio potest. nos autem id videamus quod in nobis ipsis esse debet, ut quicquid accideret fortiter et sapienter feramus et accidisse hominibus meminerimus, nosque cum multum litterae tum non minimum Idus quoque Martiae consolentur.

\ciceroSection{4} suscipe nunc meam deliberationem qua sollicitor. ita multa veniunt in mentem in utramque partem. proficiscor , ut constitueram, legatus in Graeciam: caedis impendentis periculum non nihil vitare videor sed casurus in aliquam vituperationem quod rei publicae defuerim tam gravi tempore. sin autem mansero, fore me quidem video in discrimine sed accidere posse suspicor ut prodesse possim rei publicae. iam illa consilia privata sunt, quod sentio valde esse utile ad confirmationem Ciceronis me illuc venire; nec alia causa profectionis mihi ulla fuit tum cum consilium cepi legari a Caesare. tota igitur hac de re, ut soles, si quid ad me pertinere putas, cogitabis.

\ciceroSection{5} redeo nunc ad epistulam tuam. scribis enim esse rumores me ad lacum quod habeo venditurum, minusculam vero villam utique Quinto traditurum vel impenso pretio, quo introducatur, ut tibi Quintus filius dixerit, dotata Aquilia. ego vero de venditione nihil cogito nisi quid quod magis me delectet invenero. Quintus autem de emendo nihil curat hoc tempore. satis enim torquetur debitione dotis in qua mirificas Q. Egnatio gratias agit; a ducenda autem uxore sic abhorret ut libero lectulo neget esse quicquam iucundius. sed haec quoque hactenus.

\ciceroSection{6} redeo enim ad miseram seu nullam potius rem publicam. M. Antonius ad me scripsit de restitutione Sex. Clodi; quam honorifice, quod ad me attinet, ex ipsius litteris cognosces (misi enim tibi exemplum), quam dissolute, quam turpiter quamque ita perniciose ut non numquam Caesar desiderandus esse videatur facile existimabis. quae enim Caesar numquam neque fecisset neque passus esset, ea nunc ex falsis eius commentariis proferuntur. ego autem Antonio facillimum me praebui. etenim ille, quoniam semel induxit animum sibi licere quod vellet, fecisset nihilo minus me invito. itaque mearum quoque litterarum misi tibi exemplum.
